% GPQA-Diamond ARM deep semantic analysis
% Requires: \usepackage{tcolorbox}, \findingbox, \newcounter{finding} from gpqa_macro_analysis.tex
%
% Self-contained findings: every term defined in-place.
%   F1: First cognitive move locks trajectory fate
%   F2: Deduction and generic reasoning are near-perfectly anti-correlated (r=-0.908)
%   F3: Segment count U-shape risk curve
%   F4: Tool execution breaks the deduction self-loop
%   F5: Question domain sets the cognitive baseline, harness determines the deviation
%   F6: ZeroClaw is the only harness where wrong answers produce less homogeneous cognition
%   F7: On HLE, no cognitive mode discriminates correct from wrong answers

\subsection{ARM Semantic Structure: Deeper Patterns}
\label{appendix:macro-gpqa-arm-deep}

\paragraph{Definitions.}
Every model-generated response is divided into \emph{segments}, where each
segment is one assistant message in the trajectory.
Each segment is classified into exactly one of six \emph{Agentic Reasoning
Modes} (ARM) via keyword heuristics:
\begin{itemize}[leftmargin=*, itemsep=1pt]
\item \textbf{Principle-Based Deduction (PD)}: the segment derives an answer
  from scientific laws, chemical mechanisms, mathematical theorems, or
  formal computation steps.
\item \textbf{Recovery-Replanning (RR)}: the segment retracts, corrects, or
  changes direction from prior reasoning.
\item \textbf{Uncertainty Navigation (UN)}: the segment expresses hedging,
  doubt, or probabilistic language.
\item \textbf{Intuitive Commitment (IC)}: the segment asserts an answer
  directly without extended derivation.
\item \textbf{Systematic Elimination (SE)}: the segment compares options
  and rules them out one by one.
\item \textbf{Undifferentiated reasoning (reason)}: the segment contains
  deliberation that does not trigger any of the above classifiers.
\end{itemize}

\emph{Arm entropy} measures how evenly a trajectory's segments are spread
across modes. A value near zero means nearly all segments belong to one
mode; higher values mean the trajectory uses several modes in similar
proportions.

Three \emph{harness architectures} are evaluated on the GPQA-Diamond
benchmark of 198 expert-level science multiple-choice questions, all using
the Qwen3.5-27B model.
\textbf{OpenClaw} provides an open-ended loop: the model can call tools
(shell, search, read) or emit an answer at any time, with no structural
commit deadline.
\textbf{OpenCode} provides a code-execution environment: the model can write
and run Python code, with each tool result anchoring subsequent reasoning.
\textbf{ZeroClaw} imposes a fixed template: the model must plan, then
reason for a fixed number of steps, then answer, leaving no room to re-enter
earlier phases.

% ===========================================================================
\paragraph{The first cognitive move predicts the trajectory's eventual
cognitive profile.}
% ===========================================================================

Table~\ref{tab:arm-first-mode} groups trajectories by whether their very
first segment is classified as deduction (PD) or undifferentiated reasoning
(reason).
All other starting modes (IC, UN, RR, SE) account for under 3\% of
trajectories and are omitted.

\begin{table}[t]
\centering
\small
\caption{%
  How the first cognitive move shapes the rest of the trajectory, on
  GPQA-Diamond.
  ``First mode'' = the ARM label of the very first assistant-generated
  segment.
  ``Mean PD rate'' = what fraction of all subsequent segments in the
  trajectory are classified as principle-based deduction, regardless of
  how the trajectory started.
  ``Acc'' = percentage of trajectories in this group that answer correctly.
}
\label{tab:arm-first-mode}
\setlength{\tabcolsep}{3.5pt}
\begin{tabular}{llrrrrr}
\toprule
Harness & First mode & $n$ & Acc\,\% & Mean PD rate & Mean segments & Arm ent. \\
\midrule
\multirow{2}{*}{OpenClaw}
  & PD     &  97 & 61 & \textbf{0.83} &  3.7 & 0.25 \\
  & reason &  95 & 69 & 0.30          &  4.6 & 0.70 \\
\midrule
\multirow{2}{*}{OpenCode}
  & PD     &  99 & 71 & \textbf{0.55} &  9.8 & 0.45 \\
  & reason &  93 & 74 & 0.24          & 11.4 & 0.68 \\
\midrule
\multirow{2}{*}{ZeroClaw}
  & PD     &  99 & 72 & \textbf{0.45} & 16.2 & 0.48 \\
  & reason &  93 & 84 & 0.18          & 16.4 & 0.79 \\
\bottomrule
\end{tabular}
\end{table}

\findingbox{The first cognitive move a model makes is self-reinforcing.
Trajectories that begin with deduction end up spending 2.5 to 3.5 times
more of their total cognitive budget on deduction than trajectories that
begin with generic reasoning, regardless of harness architecture.}

The effect is strongest under OpenClaw, where an initial deduction segment
leads to a mean PD rate of 0.83: the model spends over four-fifths of its
cognitive effort re-deriving principles.
Trajectories that begin with undifferentiated reasoning keep their PD rate
to 0.30, leaving room for commitment, recovery, and mode switching.

The accuracy gap consistently favors starting with reason: 6 percentage
points under OpenClaw and OpenCode, 12 points under ZeroClaw.
The mechanism is the PD self-loop: once the model enters deduction, the
probability of staying in deduction on the next segment rises sharply, and
the exit path toward commitment narrows.
A trajectory that begins with reason can still enter deduction later, but
the reverse rarely happens, because deduction is an absorbing cognitive
state.

ZeroClaw's fixed template controls the number of segments (16.2 vs.\ 16.4
regardless of first mode) but cannot control their cognitive content.
The PD rate gap (0.45 vs.\ 0.18) persists, and the accuracy gap is the
widest of the three harnesses, suggesting that when the template forces a
commit deadline, entering through deduction rather than reasoning is
particularly harmful.

% ===========================================================================
\paragraph{Deduction and generic reasoning are near-perfect substitutes on
hard-science questions.}
% ===========================================================================

On GPQA-Diamond, the Pearson correlation between a trajectory's PD rate and
its reason rate is $r = -0.908$ across all 586 agent trajectories.
No other mode pair has $|r| > 0.2$.
On the HLE benchmark (323 to 591 text-based multiple-choice questions across
dozens of academic disciplines), the PD-reason correlation is $r = -0.21$,
and the strongest anti-correlation is between uncertainty and recovery at
$r = -0.37$.

\findingbox{On hard-science benchmarks, cognitive mode allocation reduces to
a single dimension: every segment is either deduction or generic reasoning,
and the two are near-perfect substitutes.
On knowledge-intensive benchmarks, the mode space is genuinely
multidimensional, with recovery, uncertainty, and reasoning operating as
independent cognitive strategies.}

This structural fact explains why the PD rate carries such strong diagnostic
weight on GPQA-Diamond but none on HLE.
In a near-one-dimensional mode space, an increase in deduction necessarily
displaces generic reasoning.
The exit path from deduction (reason$\to$IC, the transition from generic
reasoning to commitment) collapses because there is no generic reasoning
left to serve as the exit vehicle.
On HLE, deduction and reasoning are independent dimensions: a model can
increase recovery or uncertainty without reducing reasoning, and no single
mode dominates the diagnostic landscape.

% ===========================================================================
\paragraph{Too few segments predict failure as reliably as too many.}
% ===========================================================================

Table~\ref{tab:arm-segment-bins} groups trajectories by how many segments
they contain and reports accuracy within each size bin.
A \emph{segment} is one assistant message, classified into one ARM mode.

\begin{table}[t]
\centering
\small
\caption{%
  Accuracy and cognitive profile by segment count, on GPQA-Diamond.
  ``Segments'' = number of assistant-generated messages in the trajectory.
  ``PD rate'' = fraction of segments classified as principle-based deduction.
  ``Arm ent.'' = cognitive mode diversity (higher = more modes used).
}
\label{tab:arm-segment-bins}
\setlength{\tabcolsep}{3pt}
\begin{tabular}{llrrrrr}
\toprule
Harness & Segments & $n$ & Acc\,\% & PD rate & Arm ent. & Tokens \\
\midrule
\multirow{5}{*}{OpenClaw}
  & 1--2   & 112 &  52 & 0.70 & 0.32 & 18{,}000 \\
  & 3--5   &  38 &  87 & 0.46 & 0.67 &  4{,}000 \\
  & 6--10  &  40 &  88 & 0.33 & 0.68 &  1{,}400 \\
  & 11--20 &   6 &  67 & 0.26 & 0.85 &  2{,}900 \\
\midrule
\multirow{5}{*}{OpenCode}
  & 1--2   &  23 &  \textbf{0} & 0.63 & 0.42 &  1{,}400 \\
  & 3--5   &  70 &  73 & 0.60 & 0.62 &  1{,}200 \\
  & 6--10  &  60 &  93 & 0.26 & 0.62 &  1{,}200 \\
  & 11--20 &  24 &  83 & 0.18 & 0.61 &  4{,}200 \\
  & 20+    &  21 &  86 & 0.06 & 0.32 & 10{,}000 \\
\midrule
\multirow{2}{*}{ZeroClaw}
  & 1--10  &   0 & -- & --  & --  & -- \\
  & 11--20 & 190 &  81 & 0.30 & 0.68 &  2{,}100 \\
\bottomrule
\end{tabular}
\end{table}

\findingbox{Segment count exhibits a U-shaped relationship with accuracy:
trajectories with only 1 or 2 segments are near-certain failures under
OpenCode (0\% accuracy), while the optimal range is 3 to 10 segments (73\%
to 93\% accuracy), where the model deliberates without either committing
prematurely or spiraling indefinitely.}

Under OpenCode, the 23 trajectories with 1 or 2 segments form a pure failure
class: the model emits an answer after minimal deliberation, and none are
correct.
These are premature commitments, not deduction spirals.
The 3 to 10 segment range is the healthy regime: the model reasons, invokes
code, verifies results, and converges.

Under OpenClaw, the 1 to 2 segment bin is heterogeneous.
It contains both brief correct answers (short, confident, well-formed) and
giant deduction monoliths: the mean token count is 18,000, compared to 1,400
for 6 to 10 segment trajectories.
These are single-segment PD spirals that happen to be classified as one
segment because the model never transitions to another mode before hitting
the context limit.
The 3 to 10 segment range achieves 87\% to 88\% accuracy under OpenClaw as
well, confirming that multi-step structure is necessary for reliable
correctness.

ZeroClaw's fixed template produces a narrow segment distribution centered at
17 segments regardless of outcome, so segment count carries no diagnostic
content for this harness.

% ===========================================================================
\paragraph{Executing a tool breaks the deduction self-loop.}
% ===========================================================================

Table~\ref{tab:arm-tool-catalyst} compares OpenCode trajectories that call
at least one code-execution tool against those that never do.
A \emph{transition} is a consecutive pair of segments: PD$\to$PD means a
deduction segment is immediately followed by another deduction segment.
Reason$\to$IC means a generic reasoning segment is immediately followed by
an intuitive commitment (a direct answer).

\begin{table}[t]
\centering
\small
\caption{%
  How tool use changes cognitive transitions, for OpenCode on GPQA-Diamond.
  ``Has tool use'' = the trajectory contains at least one code-execution step.
  PD$\to$PD = probability that a deduction segment is followed immediately
  by another deduction segment.
  Reason$\to$IC = probability that generic reasoning is followed immediately
  by a direct commitment to an answer.
}
\label{tab:arm-tool-catalyst}
\setlength{\tabcolsep}{4pt}
\begin{tabular}{lrrrrr}
\toprule
Group & $n$ & PD$\to$PD & Reason$\to$PD & Reason$\to$IC & Acc\,\% \\
\midrule
Has tool use    &  88 & \textbf{0.14} & 0.20 & 0.14 & \textbf{91} \\
No tool use     & 110 & \textbf{0.44} & 0.46 & 0.47 & 60 \\
\bottomrule
\end{tabular}
\end{table}

\findingbox{Executing code functions as a cognitive circuit breaker.
OpenCode trajectories that invoke tools have a deduction-to-deduction
transition probability of 0.14, compared to 0.44 for those that never invoke
tools: a threefold reduction in the self-loop that characterizes deduction
spirals.}

The mechanism is structural, not informational.
When the model executes code, the tool result arrives as a separate event
that must be processed.
The model's next segment after receiving a tool result is classified as
generic reasoning (reflection on the output), not as deduction.
This mandatory mode boundary interrupts what would otherwise be a continuous
deduction chain.

The accuracy gap is large (91\% vs.\ 60\%), but this does not mean tool use
causes correctness.
Task structure determines both whether a computational path exists and
whether the problem is solvable.
Tasks that afford tool use tend to be decomposable (numerical, algebraic)
and therefore tractable; tasks that do not afford tool use tend to require
irreducible verbal reasoning about molecular structure or biological
mechanisms, where the model defaults to deduction without a computational
anchor.
The causality runs from task structure through tool availability to both the
cognitive transition pattern and the outcome, not from the tool call itself.

% ===========================================================================
\paragraph{Question domain sets the cognitive baseline; harness architecture
determines whether that baseline helps or hurts.}
% ===========================================================================

Table~\ref{tab:arm-subdomain} reports ARM mode rates by high-level scientific
domain, pooled across all three agent harnesses.

\begin{table}[t]
\centering
\small
\caption{%
  Cognitive mode allocation by scientific domain on GPQA-Diamond, pooled
  across OpenClaw, OpenCode, and ZeroClaw.
  PD rate = fraction of segments classified as principle-based deduction.
  RR rate = fraction classified as recovery (retraction or correction).
  UN rate = fraction classified as uncertainty expression.
}
\label{tab:arm-subdomain}
\setlength{\tabcolsep}{4pt}
\begin{tabular}{lrrrrrr}
\toprule
Domain & $n$ & PD rate & IC rate & RR rate & UN rate & Acc\,\% \\
\midrule
Physics   & 258 & 0.31 & 0.04 & 0.02 & 0.01 & 89 \\
Chemistry & 279 & \textbf{0.53} & 0.01 & 0.01 & 0.02 & 62 \\
Biology   &  57 & 0.40 & 0.03 & 0.02 & 0.05 & 60 \\
\bottomrule
\end{tabular}
\end{table}

\findingbox{The cognitive strategy a model adopts is primarily determined by
the scientific domain of the question, not by the harness that executes it.
Chemistry questions demand 1.7 times more deduction than Physics questions
regardless of which harness is used.
Harness architecture determines whether that elevated deduction baseline
produces a pathological spiral or a structured derivation within a bounded
budget.}

Physics tasks (quantum mechanics, high-energy physics, electromagnetism)
have sequential algebraic structure: the model can compute a quantity,
verify the result, and commit.
Chemistry tasks (organic synthesis, stereochemistry, reaction mechanisms)
require multi-step verbal reasoning about three-dimensional molecular
structure: there is no computational shortcut, so the model defaults to
deduction and stays there.

The PD rate gap between Chemistry (0.53) and Physics (0.31) is larger than
the outcome-conditioned PD gap within any single harness.
Domain determines the cognitive baseline; harness determines the deviation
from it.
Physics tasks are decomposable and executable: they work well under every
harness (89\% pooled accuracy).
Chemistry tasks are irreducibly verbal: under OpenClaw's open loop the
baseline PD amplifies into a spiral, while under ZeroClaw's fixed template
it merely produces an elevated PD rate within a bounded budget.

% ===========================================================================
\paragraph{ZeroClaw is the only harness where wrong answers produce
\emph{less} homogeneous cognition than correct ones.}
% ===========================================================================

Table~\ref{tab:arm-dominance} reports the \emph{dominant rate}: the fraction
of a trajectory's segments assigned to its single most frequent cognitive
mode.
A dominant rate of 1.0 means every segment in the trajectory has the same
ARM label; a lower value means the trajectory mixes several modes.

\begin{table}[t]
\centering
\small
\caption{%
  Cognitive homogeneity by harness and outcome on GPQA-Diamond.
  ``Dominant rate'' = fraction of segments assigned to the trajectory's most
  frequent mode (higher = more homogeneous).
  ``Arm ent.'' = cognitive mode diversity (higher = modes used more evenly).
  ``Dominant mode'' = the most frequent mode, with the percentage of
  trajectories for which it is dominant.
}
\label{tab:arm-dominance}
\setlength{\tabcolsep}{3.5pt}
\begin{tabular}{llrrrl}
\toprule
Harness & Outcome & Dom. rate & Arm ent. & Dominant mode & \\
\midrule
\multirow{2}{*}{OpenClaw}
  & C & 0.70 & 0.52 & reason (55\% of traj.) \\
  & W & \textbf{0.74} & 0.40 & PD (60\% of traj.)    \\
\midrule
\multirow{2}{*}{OpenCode}
  & C & 0.70 & 0.60 & reason (67\% of traj.) \\
  & W & \textbf{0.73} & 0.46 & PD (55\% of traj.)    \\
\midrule
\multirow{2}{*}{ZeroClaw}
  & C & 0.71 & 0.64 & reason (77\% of traj.) \\
  & W & \textbf{0.62} & 0.81 & reason (68\% of traj.) \\
\bottomrule
\end{tabular}
\end{table}

\findingbox{ZeroClaw is the only harness where wrong trajectories are
\emph{less} cognitively homogeneous than correct ones.
Under OpenClaw and OpenCode, wrong answers coincide with mode collapse (the
model fixates on a single cognitive strategy, typically deduction).
Under ZeroClaw, wrong answers coincide with mode oscillation (the model
switches between strategies without converging on any of them).}

Under OpenClaw and OpenCode, the dominant rate rises on wrong trajectories
(from 0.70 to 0.74 and 0.73 respectively), and the dominant mode shifts
from reason (correct) to deduction (wrong).
The model collapses toward pure deduction.

Under ZeroClaw, the dominant rate \emph{drops} on wrong trajectories (from
0.71 to 0.62), and arm entropy rises from 0.64 to 0.81.
The fixed template forces the model to produce a set number of reasoning
steps.
When the model is uncertain, it fills those mandatory steps with diverse,
oscillating modes (deduction, recovery, uncertainty, reasoning) rather than
converging.
Reason remains the dominant mode for both outcomes under ZeroClaw; the
template prevents the shift to deduction-dominance that occurs under the
other two harnesses, but it cannot prevent the uncertainty that produces
mode oscillation.

% ===========================================================================
\paragraph{On knowledge-intensive benchmarks, cognitive strategy does not
predict correctness.}
% ===========================================================================

Table~\ref{tab:arm-hle-nondiscriminative} reports ARM mode rates on the HLE
benchmark, which consists of 323 to 591 text-based multiple-choice questions
spanning philosophy, linguistics, theoretical mathematics, experimental
science, and other academic disciplines.
DirectLLM is the model without any harness scaffolding (single-pass
inference); OpenCode provides a code-execution environment.

\begin{table}[t]
\centering
\small
\caption{%
  Cognitive mode rates by outcome on the HLE benchmark.
  $\Delta$ = wrong minus correct, in percentage points.
  A positive $\Delta$ means the mode is more common on wrong trajectories.
  ``$n$ seg.'' = mean number of assistant segments per trajectory.
}
\label{tab:arm-hle-nondiscriminative}
\setlength{\tabcolsep}{3pt}
\begin{tabular}{llrrrrrrr}
\toprule
Harness & Out. & PD & RR & UN & IC & Arm ent. & $n$ seg. \\
\midrule
\multirow{3}{*}{DirectLLM}
  & C      & 0.13 & 0.27 & 0.10 & 0.08 & 1.23 & 44 \\
  & W      & 0.11 & 0.26 & \textbf{0.14} & 0.07 & 1.14 & \textbf{52} \\
  & $\Delta$ & $-$2 & $-$1  & \textbf{+4}  & $-$1  & $-$9  & \textbf{+8} \\
\midrule
\multirow{3}{*}{OpenCode}
  & C      & 0.16 & 0.36 & 0.20 & 0.13 & 0.17 & 3 \\
  & W      & 0.14 & 0.31 & 0.24 & 0.15 & 0.13 & 4 \\
  & $\Delta$ & $-$2 & $-$5  & +4    & +2    & $-$4  & +1 \\
\bottomrule
\end{tabular}
\end{table}

\findingbox{On HLE, no cognitive mode rate meaningfully discriminates
correct from wrong trajectories.
The largest outcome-conditioned gap is 4 percentage points (uncertainty
expression, slightly higher on wrong answers).
This is an order of magnitude smaller than the deduction gap on GPQA-Diamond
(7 to 15 points), and it runs in the opposite direction.
Cognitive strategy describes how the model thinks; it does not predict
whether the model succeeds.}

On GPQA-Diamond, the model's cognitive strategy (how much deduction, whether
it can exit to commitment) is a cause of success or failure.
On HLE, cognitive strategy is a response to task difficulty, not a
determinant of outcome: the model explores, recovers, and hedges at nearly
identical rates whether it ultimately answers correctly or not.
What separates success from failure is the depth of domain knowledge the
model can bring to bear, which manifests in trajectory length (44 vs.\ 52
segments for DirectLLM) rather than in mode allocation.

OpenCode's arm entropy on HLE is extremely low (0.13 to 0.17) because its
compressed reasoning budget allows only 3 to 4 segments per trajectory.
In that narrow window, mode diversity is structurally impossible.
DirectLLM, unconstrained, produces 44 to 52 segments with arm entropy above
1.1, reflecting genuine multi-mode exploration.
Yet the mode distribution is nearly identical across outcomes for both
harnesses, confirming that ARM modes on HLE describe what the model chooses
to do, not whether it succeeds.

% ===========================================================================
\paragraph{What is common across both benchmarks?}
% ===========================================================================

The preceding findings treat GPQA-Diamond and HLE separately.
Table~\ref{tab:arm-cross-direction} pools all 1,500 trajectories (586 GPQA
+ 914 HLE) and reports, for each cognitive metric, whether the direction of
the wrong-minus-correct gap is the same on both benchmarks.

\begin{table}[t]
\centering
\small
\caption{%
  Direction of the wrong-minus-correct gap ($\Delta$) for each cognitive
  metric, computed separately on GPQA-Diamond and HLE and then compared.
  ``Same direction'' means the metric shifts the same way on wrong
  trajectories regardless of benchmark.
  ``Flips direction'' means the metric rises on wrong trajectories for one
  benchmark and falls on the other.
}
\label{tab:arm-cross-direction}
\setlength{\tabcolsep}{4pt}
\begin{tabular}{lrrrr}
\toprule
Metric & GPQA $\Delta$ & HLE $\Delta$ & Same? & Interpretation \\
\midrule
Arm entropy (mode diversity)    & $-$0.07 & $-$0.22 & yes & Wrong = less diverse \\
Uncertainty expression (UN)     & +0.03   & +0.05   & yes & Wrong = more hedging \\
Generic reasoning (reason)      & $-$0.16 & $-$0.03 & yes & Wrong = less deliberation \\
Number of segments              & $-$4.2  & $-$2.7  & yes & Wrong = fewer steps \\
Dominant mode purity            & +0.00   & +0.09   & yes & Wrong = more homogeneous \\
\midrule
Principle-based deduction (PD)  & \textbf{+0.14} & \textbf{$-$0.01} & \textbf{no} & GPQA: wrong = more deduction \\
                                &         &         &      & HLE: wrong = slightly less \\
Intuitive commitment (IC)       & $-$0.02 & +0.02   & no   & Flips weakly \\
Recovery-replanning (RR)        & $-$0.00 & $-$0.03 & yes  & Wrong = slightly less recovery \\
\bottomrule
\end{tabular}
\end{table}

\findingbox{Five of eight cognitive metrics shift in the same direction on
wrong trajectories regardless of whether the benchmark is hard-science
(GPQA) or knowledge-intensive (HLE).
Wrong trajectories are less cognitively diverse, more uncertain, contain
less generic deliberation, take fewer steps, and are more homogeneous.
The only metric that flips direction is deduction rate, confirming that the
PD sink is a benchmark-specific pathology, not a universal cognitive law.}

The five universal markers describe a common failure profile: the model's
cognition narrows (lower arm entropy, higher dominant rate, fewer segments)
while its expressed uncertainty rises.
This is visible even on HLE, where the overall ARM profile is
outcome-invariant at the mode-rate level: the \emph{direction} of each
metric's shift is the same, even though the \emph{magnitude} is smaller
than on GPQA.

The only genuine reversal is deduction rate.
On GPQA, wrong trajectories spend 14 percentage points more of their
cognitive budget on deduction; on HLE, wrong trajectories spend 1 point
less.
The deduction spiral (Finding~1 and Finding~3 of the main GPQA analysis) is
therefore not a property of the model, nor of agent harnesses in general,
but of the interaction between a hard-science benchmark and an open-loop
harness architecture.

\findingbox{Correct trajectories converge to approximately 70\% dominant
mode purity regardless of benchmark, harness, or absolute mode rates.
This is the only ARM metric whose central value is identical between
GPQA-correct (0.70) and HLE-correct (0.71).}

The dominant rate is the fraction of a trajectory's segments assigned to
its single most frequent cognitive mode.
A value of 1.0 means the trajectory never switches modes (pure PD, or pure
reason); a value near 0.5 means it splits evenly between two modes.

On GPQA, correct trajectories achieve this 0.70 through a mix of deduction
and generic reasoning (PD and reason are near-perfect substitutes in the
one-dimensional GPQA mode space).
On HLE, correct trajectories achieve it through a mix of recovery,
uncertainty navigation, and generic reasoning (a genuinely multidimensional
mode space).
The mechanism differs, but the outcome is the same: correct trajectories
are neither cognitively rigid (pure one mode) nor chaotic (constant mode
switching without convergence).
They settle into a dominant mode with room for occasional transitions.

Wrong trajectories deviate from this equilibrium in opposite directions
depending on harness architecture.
Under OpenClaw and OpenCode, wrong trajectories become more pure (dominant
rate 0.73 to 0.74, mode collapse toward deduction).
Under ZeroClaw, wrong trajectories become less pure (dominant rate 0.62,
mode oscillation between deduction, recovery, and uncertainty).
Both deviations from the 0.70 equilibrium predict failure.
